\section{The complete-intersection family and its bundles}
\label{sec:geometry}

We use the polarized conifold family of \cite[Section 2]{W_Hodge}.
The construction below retains its period coordinates and its exact
polarization. In particular, the bundle parameters will be independent
of the parameters which smooth the nodes.

\subsection{The compact family and the metric input}

Put $\Sigma=\PP^1_f\times\PP^1_g$ and let
\[
 Z=\PP_\Sigma(\OO\oplus\OO(1,-1))
       \times_\Sigma\PP_\Sigma(\OO\oplus\OO(1,-1)),
 \qquad k_i=g+\xi_i,
\]
where the projective bundles parametrize one-dimensional quotients.
A general anticanonical hypersurface $Y\subset Z$ has a small
contraction $Y\to X_0$ with eight ordinary double points. Four
exceptional curves have each of the two numerical types
\[
 (1,0,0,0),\qquad(0,1,-1,-1)
\]
in the divisor basis $(f,g,\xi_1,\xi_2)$. Write $I_f,I_g$ for the
corresponding groups of nodes. The compact construction and its
intersection data are given in \cite[Appendix A]{W_Toric}; the
polarized smoothing is developed in \cite[Section 8]{W_HodgeRay}
and \cite[Section 2]{W_Hodge}.

Fix the resulting proper polarized family $\pi:\cX\to\mathcal B$,
where $\mathcal B$ is a ball about zero in
\begin{equation}\label{eq:period-base}
 \left\{p\in\CC^8:\ \sum_{i\in I_f}p_i=0,
                         \quad\sum_{i\in I_g}p_i=0\right\}.
\end{equation}
Its relative polarization is denoted by $\mathcal H$. The relative
canonical bundle is trivial, and the local node equations are
\begin{equation}\label{eq:node-functions}
 z_1^2+z_2^2+z_3^2+z_4^2=\tau_i(p),\qquad
 \tau_i(p)=p_i/u_i(p),
\end{equation}
where each $u_i$ is a nonvanishing holomorphic function. These
identities are exact. The smooth locus is
$\mathcal B^\circ=\{p\in\mathcal B:\prod_i p_i\ne0\}$.
After shrinking the ball, all the units and their inverses are
bounded. We permit every path $p\to0$ in $\mathcal B^\circ$;
no ratio or phase condition is imposed on the $\tau_i$.

Let $\omega_p$ be the actual Ricci-flat Kähler form in
$c_1(\mathcal H_p)$ and let $g_p$ be its Riemannian metric.
We recall explicitly the analytic input from
\cite[Theorem 9.1 and Section 10]{W_Hodge}. There are conic radii
$\varrho_p$, equal up to fixed constants to
$(|z|^2+|\tau_i|)^{1/3}$ on the node charts and to one on a fixed
exterior, and initial Kähler metrics $\widetilde g_p^R$ which are
normalized Stenzel metrics on the inner charts. For fixed small
$R>0$ and every fixed $k$, there are $\mu>\eta>0$ such that
\begin{equation}\label{eq:ricci-flat-input}
 \limsup_{\substack{p\to0\\p\in\mathcal B^\circ}}
 \max_{j\le k}\sup_{X_p}\varrho_p^j
 |\nabla^j(g_p-\widetilde g_p^R)|_{\widetilde g_p^R}
       \le C_k R^{\mu-\eta}.
\end{equation}
The constants allow arbitrary independent node scales. On each
fixed compact exterior, $g_p$ converges in every spatial derivative
to the conically singular central Ricci-flat metric $g_0$.
There is also a uniform positive scalar Poincaré constant.
The latter follows either from the degree-one gap in the cited
Hodge theorem, by applying $\dd$ to scalar eigenfunctions, or
from the uniform diameter estimate proved below. These inputs
concern spatial derivatives; no derivative in $p$ is used.
The order in \eqref{eq:ricci-flat-input} is to fix $R$, take
$p\to0$, and then decrease $R$.

\subsection{A relative complete intersection}

Let $\mathbf P=\PP^3\times\PP^3$. The globally generated classes
$k_1,k_2$ define a morphism $q_Z:Z\to\mathbf P$.
Writing $[s_0:s_1],[t_0:t_1]$ on $\Sigma$, its local expression is
\[
 (x;y)=(v_1s;u_1t),\qquad
 (x';y')=(v_2s;u_2t).
\]
The quotient coordinates $(u_i,v_i)$ transform with the indicated
line bundles. The image is the fourfold
\begin{equation}\label{eq:ambient-central}
 V_0=\{x_0x'_1-x_1x'_0=0,\quad
                         y_0y'_1-y_1y'_0=0\}\subset\mathbf P.
\end{equation}
Indeed, each nonzero pair of $x$-vectors has a common direction
$s$, and each nonzero pair of $y$-vectors a common direction $t$.
A direction is free only if both corresponding vectors vanish;
both cannot be free. The remaining scalars recover the two
quotients. Thus the image is irreducible of dimension four.
The two equations form a regular sequence and their intersection
is generically smooth. A Cohen--Macaulay complete intersection
has no embedded components, so this identifies the reduced scheme
image as well.

The singular locus consists of the disjoint lines
\[
 \ell_x=\{x=x'=0,\ [y]=[y']\},\qquad
 \ell_y=\{y=y'=0,\ [x]=[x']\}.
\]
Near either line one determinant equation eliminates a smooth
coordinate; the other is the ordinary threefold node transverse
to the line. Hence $V_0$ is normal. The two exceptional surfaces
of $q_Z$ have codimension two and are contracted along one ruling.
Adjunction gives
\[
 K_{V_0}=\OO_{V_0}(-2,-2),\qquad
 -K_Z=2k_1+2k_2=q_Z^*\OO_{V_0}(2,2).
\]
The proper birational map onto the normal image has
$q_{Z*}\OO_Z=\OO_{V_0}$. Its anticanonical section ring is
therefore the section ring of $\OO_{V_0}(2,2)$, so $V_0$ is
exactly the anticanonical model in the compact construction.
The chosen $X_0$ is cut out in $V_0$ by a section of this bundle.
It lifts to a section $h_0$ on $\mathbf P$, because the Koszul
sequence
\[
 0\longrightarrow\OO_{\mathbf P}\longrightarrow
 \OO_{\mathbf P}(1,1)^{\oplus2}\longrightarrow
 I_{V_0}(2,2)\longrightarrow0
\]
has no higher cohomology in its first two terms.
The restriction of $h_0$ to either singular line has degree four
and four simple zeros. These are the original eight labelled nodes.

\begin{proposition}\label{prop:relative-cicy}
After shrinking $\mathcal B$, the given family admits a relative
closed embedding into $\mathbf P\times\mathcal B$ as a complete
intersection of bidegrees $(1,1),(1,1),(2,2)$. It identifies
$\mathcal H$ with $\OO_{\mathbf P}(2,2)|_{\cX}$.
The first two equations define a proper flat family
$\cV\to\mathcal B$ of normal integral Gorenstein fourfolds,
and $\cX\subset\cV$ is a relative anticanonical Cartier divisor.
\end{proposition}

\begin{proof}
The three-equation Koszul resolution gives
$H^1(X_0,\OO)=H^2(X_0,\OO)=0$. For
$M_{1,0}=\OO_{X_0}(1,0)$, the noninitial terms of that resolution
are
\[
 \OO(0,-1)^{\oplus2}\oplus\OO(-1,-2),\qquad
 \OO(-1,-2)\oplus\OO(-2,-3)^{\oplus2},\qquad
 \OO(-3,-4).
\]
They are acyclic on $\mathbf P$. Thus $h^0(M_{1,0})=4$ and
its higher cohomology vanishes; the same is true of
$M_{2,0}=\OO_{X_0}(0,1)$.

The family $\cX$ is the analytic pullback of a projective family
over a smooth algebraic Hilbert neighbourhood. The relative
Picard space is smooth near either class because $H^2(\OO)=0$,
and étale there because $H^1(\OO)=0$. After an étale
neighbourhood, a section through a smooth central point and a
normalized Poincaré bundle give the exact extensions
$\mathcal M_1,\mathcal M_2$; see
\cite[Theorems 2.5(2), 4.8 and Proposition 5.19]{Kleiman}.
The analytic map from $\mathcal B$ lifts by local inversion of
the étale map. Consequently neither the period coordinates nor
the functions \eqref{eq:node-functions} change.

Cohomology and base change extend the four generating sections
of each $M_{i,0}$ to frames of $\pi_*\mathcal M_i$.
Their evaluation maps remain surjective, by Nakayama's lemma
and properness. They give a map $\iota:\cX\to\mathbf P\times\mathcal B$.
The same Koszul calculation gives, with no higher cohomology,
\begin{equation}\label{eq:positive-twist-dimensions}
 h^0(\OO_{X_0}(1,1))=16-2=14,\qquad
 h^0(\OO_{X_0}(2,2))=100-32-1+1=68.
\end{equation}
For $(1,1)$ only the terms $\OO(1,1)$ and $\OO^{\oplus2}$
have ambient cohomology. For $(2,2)$ they are
$\OO(2,2)$, $\OO(1,1)^{\oplus2}\oplus\OO$, and $\OO$.
Their degree-zero complex is exact in negative total degree,
which proves the assertions rather than only their Euler characteristics.
The sixteen product sections span the first space in
\eqref{eq:positive-twist-dimensions}. This spanning persists
near zero, as does very ampleness of $\mathcal M_1\otimes\mathcal M_2$.
The composite of $\iota$ with the Segre embedding is consequently
a closed immersion, so $\iota$ is one as well.

If $I$ is its relative ideal, the bundles of sections of
$I(1,1)$ and $I(2,2)$ have ranks two and thirty-two. Lift the
two central determinants as a frame of the former and $h_0$
as a section of the latter. These three equations generate $I$:
the ideal is flat over the base, its restriction to $X_0$ is
the central ideal, and the cokernel of the three generators
vanishes there. Nakayama and properness again remove its support.

Work first on a neighbourhood $\mathcal T$ of the central triple
in the affine coefficient space. The regular-sequence property,
normal integrality of the fourfold, and the nonzerodivisor
property of the last equation are open there. The universal
Koszul resolutions give flatness. Pulling back by the
holomorphic coefficient map $\mathcal B\to\mathcal T$ produces
precisely $\cV$ and $\cX$, with the asserted dualizing bundles
by adjunction. Finally $\mathcal M_1^2\mathcal M_2^2$ and
$\mathcal H$ have the same central class. Their sections of the
étale relative Picard space agree near zero. Their quotient is
pulled back from the base and is trivial after normalization
along the chosen section. This proves the exact polarization
identity.
\end{proof}

\subsection{Cohomology and effective bundle deformations}

Put $L=\OO_V(1,-1)$ on any fourfold fibre of $\cV$ and set
\[
 \begin{aligned}
 E^{\rm sp}&=L\otimes\CC^m\ \oplus\ L^{-1}\otimes\CC^n,\\
 W_{m,n}&=\Hom(\CC^n,\CC^m)\oplus\Hom(\CC^m,\CC^n),\\
 G_{m,n}&=\GL_m(\CC)\times\GL_n(\CC).
 \end{aligned}
\]
The superscript denotes the split bundle. The same symbols
will denote its restriction to $X\subset V$.

\begin{proposition}\label{prop:effective-bundles}
The split family has an effective holomorphic miniversal
bundle deformation over $\mathcal B\times(W_{m,n},0)$, with
independent base and bundle parameters. Its local equivariant
germ has action
\begin{equation}\label{eq:two-arrow-action}
 (g,h)\cdot(A,B)=(gAh^{-1},hBg^{-1}).
\end{equation}
On $V$, its Ext groups in degrees zero and one are
$\mathfrak{gl}_m\oplus\mathfrak{gl}_n$ and $W_{m,n}$,
and all higher Ext groups vanish. Restriction to $X$ induces
an equivalence of formal local bundle-deformation groupoids.
On $X$ the Ext groups in degrees two and three are respectively
$W_{m,n}^*$ and $(\mathfrak{gl}_m\oplus\mathfrak{gl}_n)^*$,
but all local obstruction equations vanish.
\end{proposition}

\begin{proof}
The ambient Koszul resolution for $L^2$ is
\begin{equation}\label{eq:line-koszul}
 0\longrightarrow\OO_{\mathbf P}(0,-4)
 \longrightarrow\OO_{\mathbf P}(1,-3)^{\oplus2}
 \longrightarrow\OO_{\mathbf P}(2,-2)
 \longrightarrow\OO_V(2,-2)\longrightarrow0.
\end{equation}
Its only nonzero ambient cohomology is
$H^3(\OO_{\mathbf P}(0,-4))=\CC$, in Koszul degree $-2$.
Thus $H^1(V,L^2)=\CC$ and all other groups vanish.
Interchanging the factors proves the same assertion for
$L^{-2}$. The structure-sheaf resolution gives
$H^0(V,\OO)=\CC$ and $H^{>0}(V,\OO)=0$.
These calculations hold on the singular fibres as well.
The universal Koszul complexes on $\mathbf P\times\mathcal T$
have the same single nonzero cohomology terms, so the assertions
commute with base change. Choosing generators of
$H^3(\PP^3,\OO(-4))$ fixes algebraic frames $e_+,e_-$
for the two relative $H^1$ line bundles.

Since $E^{\rm sp}$ is locally free, its global Ext groups are
the cohomology of
\[
 F=\End E^{\rm sp}
 =\OO\otimes(\mathfrak{gl}_m\oplus\mathfrak{gl}_n)
 \oplus L^2\otimes\Hom(\CC^n,\CC^m)
 \oplus L^{-2}\otimes\Hom(\CC^m,\CC^n).
\]
This proves the ambient groups and the action
\eqref{eq:two-arrow-action}. Serre duality on the proper
Gorenstein fourfold, applied to the locally free self-dual
$F$, gives $H^i(F\otimes K_V)=H^{4-i}(F)^*$.
The sequence
\begin{equation}\label{eq:anticanonical-restriction}
 0\longrightarrow F\otimes K_V\longrightarrow F
           \longrightarrow F|_X\longrightarrow0
\end{equation}
therefore identifies Ext in degrees zero and one under
restriction and gives the stated groups in degrees two and three.

For effectivity, use the vector-bundle locus in the algebraic
stack of flat coherent sheaves on the projective flat family
$V_{\mathcal T}\to\mathcal T$. Its algebraicity, affine
finitely presented diagonal and relative deformation theory
are supplied by \cite[Tags 09DS, 0DLY, 0DLZ, 08VW and 08WA]{Stacks}.
The relative obstruction to a lift over a small Artin extension
with kernel $J$ is $\Ext^2_V(E^{\rm sp},E^{\rm sp})\otimes J=0$.
This remains the obstruction group when $V$ itself varies.
One can see it directly by lifting transition matrices on an
affine cover: the defect on triple overlaps is an
$\End E^{\rm sp}$-valued Čech two-cocycle. Its vanishing class
permits a correction, differences of lifts form the Ext-one
torsor, and infinitesimal automorphisms are Ext-zero.
Locally lifted free modules remain free.

There is an actual smooth framed atlas for this assertion.
Choose $k$ so that $E^{\rm sp}(k,k)$ is generated and has no
higher cohomology, and put $N=h^0(E^{\rm sp}(k,k))$.
The open relative Quot space of locally free quotients
$\OO_V(-k,-k)^N\to E$ inducing a frame of $H^0(E(k,k))$
carries a universal bundle. It is a $\GL_N$ framing atlas for
the bundle stack. If $K$ is the presentation kernel, the Ext
sequence gives $\Ext^1(K,E^{\rm sp})=\Ext^2(E^{\rm sp},E^{\rm sp})=0$.
The Quot obstruction theorem and finite-presentation lifting
criterion therefore make the atlas smooth over $\mathcal T$;
see \cite[Tags 09TU, 0CZU and 02H6]{Stacks}. Local freeness
on the total family follows from the flat-module criterion
\cite[Tag 00MH]{Stacks} and properness. Choose frames along
the split section. The tangent sequence there is
\[
 0\longrightarrow\Ext^0(E^{\rm sp},E^{\rm sp})
 \longrightarrow\mathfrak{gl}_N
 \longrightarrow T_{Q/\mathcal T}
 \longrightarrow\Ext^1(E^{\rm sp},E^{\rm sp})
 \longrightarrow0.
\]
A holomorphic transversal to the frame orbit through that
section, in relative smooth coordinates, has exactly the
$2mn$ parameters of $W_{m,n}$. Pullback of the universal
bundle gives the required effective family before any
formal convergence question arises.

To check restriction at every Artin order, suppose an ambient
lift has been chosen over an Artin algebra and a specified
$X$-lift exists over a small extension. First lift the ambient
bundle, using the vanishing Ext-two. The difference of its
restriction from the specified lift lies in Ext-one on $X$;
\eqref{eq:anticanonical-restriction} lifts this difference
uniquely to Ext-one on $V$. Modifying by it gives the desired
lift. For two such bundles the Hom restriction sequence has
$H^0(\Hom(-X))=H^1(\Hom(-X))=0$ at the split point.
The vanishings persist through Artin thickenings, by induction
on a small extension. Thus restriction also identifies all
Hom spaces and isomorphisms, including their prescribed
reductions. This proves the equivalence of formal groupoids.

In particular a minimal hull on $X$ has the form
$\CC[[z_1,\ldots,z_{2mn}]]/J$ with $J$ in the square of the
maximal ideal. The effective ambient family gives a map from
this hull to a power-series ring in $2mn$ variables with
invertible linear term. The induced endomorphism of the
power-series ring is an automorphism, so $J=0$.
The nonzero Ext-two group is therefore an obstruction space
with zero obstruction map. The same lifting argument,
finite presentation and semicontinuity make the restriction
morphism of bundle stacks étale near these points.

It remains to obtain the asserted equivariant germ rather
than only a transversal. On a finite-type neighbourhood,
the bundle stack is smooth over the smooth coefficient space
$\mathcal T$, has affine diagonal, and has reductive stabilizer
$G_{m,n}$ at the split point. That finite-type point is locally
closed; restrict to an open in which it is closed, as in
\cite[proof of Theorem 1.1]{AHR}. The relative étale slice
\cite[Theorem 2.1 and Remark 2.3]{dCFI} gives a commuting
étale roof to
\[
 [\,W_{m,n}\oplus T_0\mathcal T\,/G_{m,n}]
                      \longrightarrow T_0\mathcal T.
\]
The action on the second summand is trivial. The arrows toward
the normal quotients are strongly étale; the arrow toward
the bundle stack can be affine and hence representable.
All these are effective algebraic charts.

The given split section defines $\mathcal T\times BG_{m,n}$
in the bundle stack. Pull back the representable étale chart
along it. The branch through the marked point is analytically
invertible after shrinking $\mathcal T$. The connected group
$G_{m,n}$ fixes this discrete étale branch, so it descends
and lifts the entire split section. Trivialize its torsor
locally. Strong étaleness preserves its full stabilizer;
its $W_{m,n}$-coordinate is therefore fixed by $G_{m,n}$.
But $W_{m,n}^{G_{m,n}}=0$, because $(tI_m,I_n)$ has weights
$1,-1$ on the two summands. Thus the lifted section is exactly
the zero section. Base change of the strongly étale arrow
to $W_{m,n}\oplus T_0\mathcal T$ and local inversion give a
$G_{m,n}$-equivariant inverse germ, by uniqueness of the
branch. Pull back its actual bundle and then the fixed map
$\mathcal B\to\mathcal T$. Restriction to $X$ preserves
this germ by the étale equivalence just proved.
\end{proof}

Equivariance here is a statement about germs: for a fixed
$g\in G_{m,n}$ the equality holds wherever its two local
representatives are defined. It does not require a bounded
parameter ball invariant under the whole complex group.
We shall choose a common small ball invariant under the
compact group $U(m)\times U(n)$.

\subsection{The rank-four family and its characteristic classes}

From now on $m=n=2$. Put $D=c_1(L|_X)$ and
$H=c_1(\mathcal H|_X)$. If $a,b$ are the two ambient hyperplane
classes, then $[X]=2(a+b)^3$ in $\mathbf P$ and
$\int_{\mathbf P}a^3b^3=1$. Consequently
\[
 a^3=b^3=2,\qquad a^2b=ab^2=6,
 \qquad D=a-b,\quad H=2a+2b,
\]
and
\begin{equation}\label{eq:chern-numbers}
 H^3=320,\qquad DH^2=0,\qquad D^2H=-16,
 \qquad V:=\Vol(X_p,g_p)=\frac{160}{3}.
\end{equation}
The rank-four family $\cE_{p,z}$, $z=(A,B)\in
W:=\End(\CC^2)\oplus\End(\CC^2)$, has
\begin{equation}\label{eq:bundle-chern}
 \det\cE_{p,z}=\OO,\qquad c_2(\cE_{p,z})=-2D^2,
 \qquad c_3(\cE_{p,z})=0,\qquad c_2(\cE_{p,z})H=32.
\end{equation}
The Chern identities follow from $(1+D)^2(1-D)^2$ and
constancy in the deformation family. For the determinant
itself, $H^1(V,\OO)=0$ makes the relative Picard discrete
near the trivial bundle. Its class stays trivial and its
remaining base line bundle is trivialized locally.
This supplies a determinant trivialization for the family.

The lines $L,L^{-1}$ have zero slope and are distinct,
since $D^2H\ne0$. Thus $E^{\rm sp}_p$ is slope polystable
with these two stable types and multiplicity two each.
Interchanging the ambient factors in the Koszul Euler
characteristic shows that the two lines have equal Hilbert
polynomials. The split bundle is also Gieseker polystable,
on the central fibre as well as nearby fibres. The nonzero
pairing in \eqref{eq:bundle-chern} excludes a flat connection.
Every sufficiently small matrix parameter describes an
actual integrable bundle; its stability and its
Hermitian--Einstein metrics will be determined below.
